\documentclass[12pt]{article}
\usepackage{tcolorbox}
\usepackage{amsmath}

\include{preamble}
%All credit goes towards Dr. Adam Kapelner for providing the preamble and homework template.

\newtoggle{professormode}




\title{MATH 241 Spring 2026 Homework \#6}

\author{Steven Ou} %STUDENTS: write your name here

\iftoggle{professormode}{
\date{Due via Brightspace by 11:59PM Sunday, April $12^{nd}$, 2026 \\ \vspace{0.5cm} \small (this document last updated \today ~at \currenttime)}
}

\renewcommand{\abstractname}{Instructions and Philosophy}

\begin{document}
\maketitle

\iftoggle{professormode}{
\begin{abstract}
The path to success in this class is to do many problems. Unlike other courses, exclusively doing reading(s) will not help. Coming to lecture is akin to watching workout videos; thinking about and solving problems on your own is the actual ``working out.''  Feel free to \qu{work out} with others; \textbf{I want you to work on this in groups.}

The problems below are color coded: \ingreen{green} problems are considered \textit{easy} and marked \qu{[easy]}; \inorange{yellow} problems are considered \textit{intermediate} and marked \qu{[harder]}, \inred{red} problems are considered \textit{difficult} and marked \qu{[difficult]} and \inpurple{purple} problems are extra credit. The \textit{easy} problems are intended to be ``giveaways'' if you went to class. Do as much as you can of the others; I expect you to at least attempt the \textit{difficult} problems. 

Bonus points are given as a bonus if the homework is typed using \LaTeX. Links to installing \LaTeX~and programs for compiling \LaTeX~is found on the syllabus. You are encouraged to use \url{https://www.overleaf.com/}. If you are handing in homework this way, read the comments in the code; there is a line to comment out and you should replace my name with yours. The easiest way to use Overleaf is to go to \textit{Menu} and select \textit{Copy Project}. If you are asked to make drawings, you can take a picture of your handwritten drawing and insert them as figures or leave space using the \qu{$\backslash$vspace} command and draw them in after printing or attach them stapled.


\end{abstract}

\thispagestyle{empty}
\vspace{1cm}
NAME: \line(1,0){380}
\clearpage
}

\problem{In this exercise, we will review the beginning of the material pertaining to Section 3.1.}



\begin{enumerate}

\inthenotessubproblem{Formally define what is meant by the conditional probability of the event $A$ given that the event $B$ has occurred.}
\\
The conditional probability of event A given that event B has occurred is denoted by P(A|B). If P(B) >0, it is defined as: 
P(A|B) = $\frac{P(A\cap B)}{P(B)}$

\inthenotessubproblem{Explain why it makes sense to define $P(A|B)$ to be $ \frac{P(A \cap B)}{P(B)}$. }
\\
When $B$ has occurred, $B$ becomes our "new sample space" or the realm of all possible outcomes. For $A$ to occur within this new space, the actual outcome must be in both $A$ and $B$ (the intersection $A \cap B$). Thus, the probability is the "stuff" in $A$ that can still happen ($A \cap B$) divided by all the "stuff" that can now happen ($B$).

\inthenotessubproblem{Consider the way that the MegaMillions (lottery) is played. Suppose we played the following numbers: the white balls selected will be 20, 30, 54, 63, 65 and the golden number will be 14. Let $A = \{ \text{we win the lottery} \}$. Assuming that the selection of numbers is random, find $P(A)$.   }
\\The total number of outcomes $|S|$ is calculated by choosing 5 white balls from 70 and 1 golden ball from 24: $\binom{70}{5} \times \binom{24}{1} = 290,472,336$. Since we have one unique ticket,$P(A) = \frac{1}{290,472,336}$

\intermediatesubproblem{Suppose we are watching the lottery live and the first (white) number drawn is 65. What is the probability that we win the lottery ?}
\\The probability we win the lottery is $\frac{1}{\binom{69}{4} \times \binom{24}{1}}$.

\intermediatesubproblem{Suppose we are watching the lottery live and the first few (white) numbers drawn are 30, 54, 63, 65. What is the probability that we win the lottery?}
\\The probability that we win the lottery is $\frac{1}{\binom{66}{1} \times \binom{24}{1}} = \frac{1}{66 \times 24} = \frac{1}{1,584}$

\hardsubproblem{If you win the lottery, what would you do with the money?}\\
If I win the lottery, I would invest all of it into stocks markets!! Increasing the money, it doesn't hurt to have more money!

\end{enumerate}


\problem{These exercises will test some of the theoretic aspects of conditional probability.}

\begin{enumerate}


\easysubproblem{ If $S$ is the sample space of an experiment and $A$ is any event in that space, what is the value of $P(A|S)$? 
}
$$P(A|S) = \frac{P(A \cap S)}{P(S)} = \frac{P(A)}{1} = P(A)$$


\easysubproblem{Prove if $P(B)>0$, then $P(B|B) = 1$.}
$$P(B|B) = \frac{P(B \cap B)}{P(B)} = \frac{P(B)}{P(B)} = 1$$



\inthenotessubproblem{Prove if $P(B)>0$, then $P(A^c|B) = 1 - P(A|B)$}
\\First step add $- P(A|B)$ to the left side. Then:$$\frac{P(A^c \cap B) + P(A \cap B)}{P(B)} = \frac{P(B)}{P(B)} = 1$$



\inthenotessubproblem{State and prove The Basic Conditioning Rule.}\\
Statement: \\$P(A \cap B) = P(A)P(B|A)$ (if $P(A)>0$)\\ or\\ $P(A \cap B) = P(B)P(A|B)$ (if $P(B)>0$).\\Proof: Start with the RHS: \\$P(A)P(B|A) = P(A) \frac{P(B \cap A)}{P(A)} = P(B \cap A) = P(A \cap B)$.


\inthenotessubproblem{State and prove The Generalized Conditioning Rule.}\\
Statement:\\ $P(E_1 \cap E_2 \cap \dots \cap E_n) = P(E_1)P(E_2|E_1)P(E_3|E_1 \cap E_2) \dots P(E_n|E_1 \cap \dots \cap E_{n-1})$.\\
Proof:
\\$$P(E_1) \cdot \frac{P(E_1 \cap E_2)}{P(E_1)} \cdot \frac{P(E_1 \cap E_2 \cap E_3)}{P(E_1 \cap E_2)} \dots \frac{P(E_1 \cap E_2 \cap \dots \cap E_n)}{P(E_1 \cap E_2 \cap \dots \cap E_{n-1})}$$
\\After canceling them $P(E_1 \cap E_2 \cap \dots \cap E_n)$


\easysubproblem{Suppose $E, F$ and $G$ are events of a sample space $S$. How many different ways can you write $P(E \cap F \cap G)$ using The Generalized Conditioning Rule? }
\\There's 3! = 6 ways of writing it : \\$1.P(E) P(F|E) P(G|E \cap F)$\space$2.P(E) P(G|E) P(F|E \cap G)$\space
$3.P(F)P(E|F)P(G|F \cap E)$ \space$4.P(F)P(G|F)P(E|F \cap G)$ \space$5.P(G)P(E|G)P(F|G \cap E)$ \space$6.P(G)P(F|G)P(E|G \cap F)$ 

\end{enumerate}
\newpage


\problem{These exercises will test some of the applications of conditional probabilities. }


\begin{enumerate}

\easysubproblem{ Two fair dice are rolled. What is the conditional probability that at least
one die lands on 6 given that the dice land on different numbers?
 }
\\ Let A = {At least one die lads on 6}\\Let B ={The dice land on different numbers}\\ B= 36 possible outcomes removing the doubles because it won't occur so 36-6=30. \\ $A\cap B =$10, results that can have 6 in it so now...\\
$P(A|B) = \frac{|A \cap B|}{|B|}$
$P(A|B) = \frac{10}{30} = \mathbf{\frac{1}{3}}$

\easysubproblem{Suppose two fair dice are rolled and let $A = \{ \text{the first die lands on 6} \}$. Let $B_i = \{ \text{the sum of the dice is}~ i \}$. Find $P(A | B_2), P(A| B_3), \ldots, P(A | B_{12})$. 
}\\
$P(A|B_2) \dots P(A|B_6) = 0$ \\(Sum cannot be $\leq 6$ if the first die is 6).\\$P(A|B_7) = 1/6$ (Outcome is $(6,1)$ out of 6 possible sums).\\$P(A|B_8) = 1/5$ (Outcome is $(6,2)$ out of 5 possible sums).\\$P(A|B_9) = 1/4$, $P(A|B_{10}) = 1/3$, $P(A|B_{11}) = 1/2$; $P(A|B_{12}) = 1/1 = 1$.


%\easysubproblem{Suppose two fair dice are rolled and let $A = \{ \text{at least one of the dice lands on 6} \}$. Let $B_i = \{ \text{the sum of the dice is}~ i \}$. Find $P(A | B_2), P(A| B_3), \ldots, P(A | B_{12})$. }

\intermediatesubproblem{You are a member of a class of 18 students. A bowl contains 18 chips: 1 blue and 17 red. Each student is to take 1 chip from the bowl without replacement. The student who draws the blue chip is guaranteed an A for the course. If you have a choice of drawing first, fifth, or last, which position would you choose? Justify your choice on the basis of probability.}
\\ The position that I would choose does not matter. Using the Generalized Conditioning Rule, the probability for the $k$-th person is:$P(k\text{-th is blue}) = P(1\text{st is red}) \dots P((k-1)\text{th is red})P(k\text{-th is blue}| \dots)$$= \frac{17}{18} \times \frac{16}{17} \times \dots \times \frac{18-k+1}{18-k+2} \times \frac{1}{18-k+1} = \frac{1}{18}$.All positions have a $1/18$ chance.
\easysubproblem{Each time a shopper purchases a tube of toothpaste,
he chooses either brand $A$ or brand $B$. Suppose that for
each purchase after the first, the probability is 1/3 that he
will choose the same brand that he chose on his preceding
purchase and the probability is 2/3 that he will switch
brands. If he is equally likely to choose either brand $A$
or brand $B$ on his first purchase, what is the probability
that both his first and second purchases will be brand $A$
and both his third and fourth purchases will be brand $B$?
}

$P(A_1) = 1/2$.$P(A_2|A_1) = 1/3$ (same brand).$P(B_3|A_2) = 2/3$ (switch).$P(B_4|B_3) = 1/3$ (same brand).$P(A_1 \cap A_2 \cap B_3 \cap B_4) = \frac{1}{2} \times \frac{1}{3} \times \frac{2}{3} \times \frac{1}{3} = \frac{2}{54} = \frac{1}{27}$.


\end{enumerate}



\problem{These questions will test the idea behind independence which is covered in Section 3.2. }


\begin{enumerate}

\inthenotessubproblem{ Formally, how do we define the word independence? }\\
Two events $E$ and $F$ are independent if $P(E \cap F) = P(E)P(F)$

\inthenotessubproblem{Intuitively, what does it mean when we say two events are independent?}\\
The occurrence of one event does not affect the probability of the other event.

\inthenotessubproblem{Explain why it is impossible for two events, $E$ and $F$, to be independent and disjoint. You may assume that $P(E) > 0$ and $P(F) > 0$.}\\
If $E$ and $F$ are independent and have $P > 0$, then $P(E \cap F) = P(E)P(F) > 0$. If they are disjoint, $P(E \cap F) = 0$. You cannot have a number be both $>0$ and $=0$ simultaneously.


\hardsubproblem{In class, we answered the rocket question via two different methods (a direct method followed by the complement method). We then considered the following problem where we used the direct method: \textit{A biased coin is repeatedly tossed until a heads appears for the first time at which point the experiment stops. Assuming that the outcomes of the tosses are independent of each other, find the probability that a heads will be obtained at some point. You may assume the probability that we get heads on any flip is 0.0001.} Try to answer this problem using the complement method.} \\
$A = \{\text{eventually get heads}\}$, \\$A^c = \{\text{never get heads / infinite tails}\}$.$P(A^c) = \lim_{n \to \infty} (0.9999)^n = 0$.\\$P(A) = 1 - P(A^c) = 1 - 0 = 1$.


\inthenotessubproblem{A sequence of five independent trials are to be performed. Each trial consists of administering a drug to a patient and then seeing if the drug is a success (beneficial) or a failure (non-beneficial) for that patient. The probability the drug is successful for any patient is 0.75. Find the probability that at least one success occurs.}
\\$P(\text{Success}) = 0.75$, \\$P(\text{Failure}) = 0.25,$\\$P(\text{at least one}) = 1 - P(\text{all failures}) = 1 - (0.25)^5 = 1 - 0.0009765625 = 0.999023$


\inthenotessubproblem{A sequence of five independent trials are to be performed. Each trial consists of administering a drug to a patient and then seeing if the drug is a success (beneficial) or a failure (non-beneficial) for that patient. The probability the drug is successful for any patient is 0.75. Find the probability that exactly four successes occur.}
\\ $\binom{5}{4} * 0.75^{4}*0.25^1=0.3955$\\The probability of exactly 4 successes occur is about 40 percent. 
\end{enumerate}

\problem{In this question, we are going to generalize the problem above.}

\begin{enumerate}
    \intermediatesubproblem{A sequence of $n$ independent trials are to be performed where each trial may result in either a success with probability $p$ or may result in a failure with probability $1-p$. Show that the probability that we obtain exactly $x$ successes (where $x$ is either $0, 1, 2, \ldots, n$) is given by $\binom{n}{x}p^x (1-p)^{n-x}$. (Please star or circle this problem! We will come back to this in a later chapter).}\\
    $P(X=x) = \binom{n}{x}p^x(1-p)^{n-x}$ \space Not going to  circle it but will give the formula.
\end{enumerate}


\end{document}
